\chapter{Conclusion}
\label{chap:conclusion}

This thesis set out to improve tool support for Object-Centric Process Mining (OCPM) by addressing recurring challenges in extensibility, maintainability, and reuse.

To capture the current state of OCPM research tools, this thesis conducted a systematic literature review and analyzed the tools described in the OCPM literature.
The results indicate a fragmented landscape in which tools either prioritize rapid implementation with little surrounding infrastructure or provide interactive applications with dedicated user interfaces.
Tools in the first direction are often delivered as collections of research code with minimal integration, which makes them harder to reuse and frequently requires manual configuration.
Tools in the second direction improve accessibility, but they often reimplement recurring boilerplate functionality, for example for data import and export, log inspection, filtering, and visualization.
Across both directions, interoperability remains limited, and many tools are difficult to set up and run in practice.
As a result, effective use often requires substantial technical knowledge.

Building on these findings, this thesis designed an extensible architecture that integrates common OCPM functionality into a shared framework.
The architecture follows a dual extension model that mirrors the two main development directions observed in existing research tools.
Plugins provide a lightweight way to package and run OCPM techniques within a unified environment.
They come with an automatically generated user interface, so developers can add functionality without implementing custom UI components.
In contrast, modules use Ocelescope as a base for building specialized applications.
They allow developers to implement specialized user interfaces and application-specific logic while reusing existing functionality, such as import and export, log inspection, and filtering.

The framework was implemented as an open-source web application to support cross-platform use.
To integrate OCPM functionality in an environment commonly used in research while supporting user interfaces built with widely used technologies, the framework combines a Python backend with a TypeScript frontend.
It includes a dedicated Python library that supports plugin development and can also be used as a general-purpose library for OCPM functionality.
The project is documented extensively and distributed via Docker to support a reproducible setup and consistent execution across platforms.
Build and release steps are automated through GitHub Actions to support consistent packaging and distribution.

The evaluation indicates that the proposed approach is feasible in practice.
End users were able to install Ocelescope and execute plugins, and plugin developers were able to implement and integrate a plugin with generally positive ratings for perceived usefulness and ease of use.
In addition, the proof-of-concept plugins demonstrate that the framework can support more demanding extension scenarios, including OCEL extensions and Rust implementations through Python bindings.

Overall, this thesis provides a practical foundation for reducing fragmentation in OCPM research tooling by enabling reuse of common functionality and structured extension through plugins and modules.

\paragraph{Future Work.}

Several directions for future work follow from this thesis.

\begin{itemize}
	\item \textbf{Improved support for plugin development.}
	      This includes clearer documentation, better guidance, and live updates during development to simplify the creation and testing of plugins.

	\item \textbf{Extended plugin functionality.}
	      Additional input types and more visualization options would make plugins suitable for a broader range of OCPM tasks.

	\item \textbf{Better user interface for running plugins.}
	      More responsive feedback, clearer progress indicators, and improved handling of long running tasks would make plugin execution easier to understand and use.

	\item \textbf{Expansion of the OCPM features in the Ocelescope package.}
	      Extending the current Ocelescope Python package with functionality such as performance analysis or variant calculations would help the framework support OCPM more completely.

	\item \textbf{Evaluation of the module system in real settings.}
	      Testing the module system on real life use cases would reveal recurring challenges.
	      The framework could then be adapted to support these scenarios more effectively.
\end{itemize}
